\title{YaRN: Efficient Context Window Extension of Large Language Models}

\begin{abstract}
Rotary Position Embeddings (RoPE) have been shown to effectively encode positional information in transformer-based language models. However, these models fail to generalize past the sequence length they were trained on. We present YaRN (Yet another RoPE extensioN method), a compute-efficient method to extend the context window of such models, requiring 10x less tokens and 2.5x less training steps than previous methods. Using YaRN, we show that LLaMA models can effectively utilize and extrapolate to context lengths much longer than their original pre-training would allow, while also surpassing previous the state-of-the-art at context window extension. In addition, we demonstrate that YaRN exhibits the capability to extrapolate beyond the limited context of a fine-tuning dataset. Code is available at \url{https://github.com/jquesnelle/yarn}.
\end{abstract}

\section{Introduction}

Transformer-based Large Language Models\citep{vaswani2017attention} (LLMs) have become the near-ubiquitous choice for many natural language processing (NLP) tasks where long-range abilities such as \emph{in-context learning} (ICL) has been crucial.

In performing the NLP tasks, the maximal length of the sequences (the \emph{context window}) determined by its training processes has been one of the major limits of a pretrained LLM. Being able to dynamically extend the context window via a small amount of fine-tuning (or without fine-tuning) has become more and more desirable. To this end, the position encodings of transformers are the center of the discussions.

The original Transformer architecture used an absolute sinusoidal position encoding, which was later improved to a learnable absolute position encoding~\citep{gehring2017convolutional}. Since then, relative positional encoding schemes~\citep{shaw2018self} have further increased the performance of Transformers. Currently, the most popular relative positional encodings are \emph{T5 Relative Bias}~\citep{roberts2019t5}, \emph{RoPE}~\citep{su2022roformer}, \emph{XPos}~\citep{sun2022lengthextrapolatable}, and \emph{ALiBi}~\citep{press2022train}.

One reoccurring limitation with positional encodings is the inability to generalize past the context window seen during training. While some methods such as ALiBi are able to do limited generalization, none are able to generalize to sequences significantly longer than their pre-trained length~\citep{kazemnejad2023impact}.

Some works have been done to overcome such limitation. \citep{chen2023extending} and concurrently~\citep{kaiokendev} proposed to extend the context length by slightly modifying RoPE via Position Interpolation (PI) and fine-tuning on a small amount of data. As an alternative, \citep{blocntkaware} proposed the "NTK-aware" interpolation by taking the loss of high frequency into account. Since then, two improvements of the "NTK-aware" interpolation have been proposed, with different emphasis:

\begin{itemize}
    \item the "Dynamic NTK" interpolation method \citep{emozillareddit} for pre-trained models without fine-tuning.
    \item the "NTK-by-parts" interpolation method \citep{blocntkparts} which performs the best when fine-tuned on a small amount of longer-context data.
\end{itemize}

The "NTK-aware" interpolation and the "Dynamic NTK" interpolation have already seen their presence in the open-source models such as Code Llama~\citep{rozière2023code} (using "NTK-aware" interpolation) and Qwen~7B~\citep{qwen} (using "Dynamic NTK").

In this paper, in addition to making a complete account of the previous unpublished works on the "NTK-aware", the "Dynamic NTK" and the "NTK-by-parts" interpolations, we present YaRN (Yet another RoPE extensioN method), an improved method to efficiently extend the context window of models trained with Rotary Position Embeddings (RoPE) including the  LLaMA~\citep{touvron2023llama}, the GPT-NeoX~\citep{black2022gptneox20b}, and the PaLM~\citep{chowdhery2022palm} families of models.

The relationship between different methods and how they evolve into YaRN can be summarized into the following diagram:

YaRN reaches state-of-the-art performances in context window extensions after fine-tuning on less than $\sim$0.1\% of the original pre-training data. In the meantime, by combining with the inference-time technique called Dynamic Scaling, the Dynamic-YaRN allows for more than 2x context window extension without any fine-tuning.

\section{Background and Related Work}

\subsection{Rotary Position Embeddings}\label{sec:rope}

The basis of our work is the Rotary Position Embedding (RoPE) introduced in~\citep{su2022roformer}. We work on a hidden layer where the set of hidden neurons are denoted by $D$. Given a sequence of vectors ${\bm{x}}_1, \cdots, {\bm{x}}_L\in \mathbb{R}^{|D|}$, following the notation of~\citep{su2022roformer}, the attention layer first converts the vectors into the query vectors and the key vectors:

\begin{align}
{\bm{q}}_m = f_q({\bm{x}}_m, m) \in \mathbb{R}^{|D|}, ~ {\bm{k}}_n = f_k({\bm{x}}_n, n) \in \mathbb{R}^{|D|}.
\end{align}

Next, the attention weights are calculated as

\begin{equation}\label{eq:attention}
\text{softmax}(\dfrac{{\bm{q}}_m^T{\bm{k}}_n}{\sqrt{|D|}}),
\end{equation}

where ${\bm{q}}_m, {\bm{k}}_n$ are considered as column vectors so that ${\bm{q}}_m^T{\bm{k}}_n$ is simply the Euclidean inner product. In RoPE, we first assume that $|D|$ is even and identify the embedding space and the hidden states as complex vector spaces:
\[\mathbb{R}^{|D|}\cong \mathbb{C}^{|D|/2}\]
where the inner product ${\bm{q}}^T{\bm{k}}$ becomes the real part of the standard Hermitian inner product $\text{Re}({\bm{q}}^*{\bm{k}})$. More specifically, the isomorphisms interleave the real part and the complex part

\begin{align}
\big(({\bm{x}}_m)_1, \cdots, ({\bm{x}}_m)_{|D|} \big) \mapsto \big(({\bm{x}}_m)_1 + i({\bm{x}}_m)_2, \cdots, (({\bm{x}}_m)_{|D|-1} + i({\bm{x}}_m)_{|D|})\big), \\
\big((\textbf{q}_m)_1, \cdots, (\textbf{q}_m)_{|D|} \big) \mapsto \big((\textbf{q}_m)_1 + i(\textbf{q}_m)_2, \cdots, ((\textbf{q}_m)_{|D|-1} + i(\textbf{q}_m)_{|D|})\big).
\end{align}

To convert embeddings ${\bm{x}}_m, {\bm{x}}_n$ into query and key vectors, we are first given $\mathbb{R}$-linear operators 
\[
{\bm{W}}_q, {\bm{W}}_k: \mathbb{R}^{|D|}\rightarrow \mathbb{R}^{|D|}.
\]
Let $\bm{\theta} = \text{diag}(\theta_1, \cdots, \theta_{|D|/2})$. In complex coordinates, we define

\begin{equation}\label{eq:rope}
f_{{\bm{W}}}({\bm{x}}_m, m, \bm{\theta}) = e^{im\bm{\theta}}{\bm{W}}{\bm{x}}_m,
\end{equation}

for any linear operator ${\bm{W}}$. The functions $f_q, f_k$ in RoPE are given by

\begin{equation}
    f_q = f_{{\bm{W}}_q}, ~ f_k = f_{{\bm{W}}_k}.
\end{equation}

where $\theta_d = b^{-2d/|D|}$ and $b=10000$. This way, RoPE associates each (complex-valued) hidden neuron with a separate frequency $\theta_d$. The benefit of doing so is that the dot product between the query vector and the key vector only depends on the relative distance $m - n$.

In later discussions, a context length interpolation usually aims to modify the equation Eq.~\ref{eq:rope}. To set up a uniform convention for these discussions, note that a modification $f'_{{\bm{W}}}$ can take the following form:

\begin{align}\label{eq:modifyrope}
    f'_{{\bm{W}}}({\bm{x}}_m, m, \bm{\theta}) = f_{{\bm{W}}}({\bm{x}}_m, g(m), {\bm{h}}(\bm{\theta})),
\end{align}

where $g(m)$ is a map between real numbers and ${\bm{h}}(\bm{\theta})$ acts on the entries of the diagonal matrix $\bm{\theta}$ uniformly by $\text{diag}(h(\theta_1), \cdots, h(\theta_{|D|/2}))$ according to a function $h$. $g$ and $h$ are method-dependent functions. 

In the subsequent sections, when we introduce a new interpolation method of the form Eq.~\ref{eq:modifyrope}, we only specify the functions $g(m)$ and $h(\theta_d)$.

\subsection{Additional notations}\label{sec:related-work}
Given the pretrained maximal context length $L$, our goal is to extend it to $L' > L$ either with or without finetuning. We introduce the notion of {\it scale factor} $s$ defined by $s = \frac{L'}{L}$.

For the convenience of some discussions, we also introduce \emph{wavelength} $\lambda_d$ associated with the $d$-th hidden dimension of RoPE as follows:

\begin{align} \label{eq:wavelength}
    \lambda_d = \dfrac{2\pi}{\theta_d} = 2\pi b^{\frac{2d}{|D|}}.
\end{align}

The wavelength describes the length of tokens needed in order for the rotary position embedding at dimension $d$ to perform a full rotation ($2\pi$).

\subsection{Related work}

Position Interpolation (PI) is one of the earlier works extending context lengths of RoPE proposed by \citet{chen2023extending}, and concurrently~\citet{kaiokendev}. Under the notation of Eq.~\ref{eq:modifyrope}, it is setting

\begin{equation}
g(m) = s\cdot m, ~{\bm{h}}(\bm{\theta}) = \bm{\theta},
\end{equation}

where $s$ is the scale factor $\frac{L'}{L}$. We include some details in Appendix \ref{appendix:pi}.

ReRoPE~\citep{rerope2023} also aims to extend the context size of existing models pre-trained with RoPE, and claims "infinite" context length without needing any fine-tuning. This claim is backed by a monotonically decreasing loss with increasing context length up to 16k on the Llama 2 13B model. It achieves context extension by modifying the attention mechanism and thus is not purely an embedding interpolation method. Since it is currently not compatible with Flash Attention~2~\citep{dao2023flashattention2} and requires two attention passes during inference, we do not consider it for comparison.

Concurrently with our work, LM-Infinite~\citep{han2023lminfinite} proposes similar ideas to YaRN, but focuses on "on-the-fly" length generalization for non-fine-tuned models. Since they also modify the attention mechanism of the models, it is not an embedding interpolation method and is not immediately compatible with Flash Attention 2.

\section{Methodology}

Whereas PI stretches all RoPE dimensions equally, we find that the theoretical interpolation bound described by PI~\citep{chen2023extending} is insufficient at predicting the complex dynamics between RoPE and the LLM's internal embeddings. In the following subsections, we describe the main issues with PI we have individually identified and solved, so as to give the readers the context, origin and justifications of each method which we use in concert to obtain the full YaRN method.

\subsection{Loss of High Frequency information - "NTK-aware" interpolation}\label{sec:ntkaware} 

If we look at rotary position embeddings (RoPE) only from an information encoding perspective, it was shown in~\citep{tancik2020fourier}, using Neural Tangent Kernel (NTK) theory, that deep neural networks have trouble learning high frequency information if the input dimension is low and the corresponding embeddings lack high frequency components. Here we can see the similarities: a token's positional information is one-dimensional, and RoPE expands it to an n-dimensional complex vector embedding. RoPE closely resembles Fourier Features~\citep{tancik2020fourier} in many aspects, as it is possible to define RoPE as a special 1D case of a Fourier Feature.

In the case of Positional Interpolation (PI), as we strech all dimensions equally by a factor $s$, it removes the high frequency components of RoPE. This degradation is worsened as the scaling factor $s$ grows, and at some point, the network will not be able to recover. Previous fine-tunes~\citep{kaiokendev}~\citep{chen2023extending}~\citep{together}~\citep{llongma} using PI were only able to achieve a scaling factor of roughly $s=8$ before the LLM's outputs starts to degrade, even after fine-tuning.

In order to alleviate this issue, the "NTK-aware" interpolation was developed in~\citep{blocntkaware}. Instead of scaling every dimension of RoPE equally by a factor $s$, we spread out the interpolation pressure across multiple dimensions by scaling high frequencies less and low frequencies more. One can obtain such a transformation in many ways, but the simplest would be to perform a base change on the value of $\theta$. The details are described in the Appendix \ref{appendix:ntkaware} and the method has seen some open-source adoptions\footnote{We note that shortly before the release of this article, Code Llama~\citep{rozière2023code} was released and uses "NTK-aware" scaling by manually scaling the base $b$ to 1M, in which they call this method as RoPE "adjusted base frequency" (ABF).}.

One main issue of this "NTK-aware" scaling is that it is very difficult to determine what optimal base should be used for an intended context extension by $s$ times. The best base to use for "NTK-aware" interpolation usually has to be found empirically, which significantly increases the difficulty and cost of obtaining a successful fine-tuned model. Despite its limitations, the observations from the NTK theory is valid and the following idea is still maintained and executed in a different way in the "NTK-by-parts" interpolation introduced in the next section.

\subsection{Loss of Relative Local Distances - "NTK-by-parts" interpolation} \label{sec:ntkparts} 

To understand why "NTK-aware" interpolation works better than PI and to eliminate its disadvantages, we have to take a closer look at RoPE. In this section, we think heavily in terms of the wavelengths $\lambda_d$ defined in Eq.~\ref{eq:wavelength} in the formula of RoPE. For simplicity, we omit the subscript $d$ in $\lambda_d$ and the reader is encouraged to think about $\lambda$ as the wavelength of an arbitrary periodic function.

In theory, as RoPE is a relative position embedding, it should be quite surprising that it fails to generalize to unseen longer context sizes. However, we can show that in practice, RoPE does not only encode relative position. One observation we can make is that given a context size $L$, there are some dimensions $d$ where the wavelength is longer than the maximum context length seen during pretraining ($\lambda > L$), this suggests that some dimensions' rotary embeddings might not be distributed evenly in the rotational domain (i.e. does not perform a full rotation for the entire training context size). In such cases, we presume having unique position pairs\footnote{Since the dimension never rotates fully at least once during pre-training, if we pick the first token as the anchor, every other token during pre-training has an unique distance to it, which the neural network can use to determine its absolute position.} implies that the absolute positional information remains intact in those dimensions. On the contrary, when the wavelength is short, only relative positional information is accessible to the network.

Given these observations, we can see that it is important to not touch the dimensions that only encode relative positional information, as they are crucial for the network to distinguish the relative order of nearby tokens. Meanwhile, dimensions that only encode absolute positional information should always be interpolated, as larger distances will be out of distribution. Instead of arbitrarily changing the base in "NTK-aware" interpolation (which basically does something similar to what is described here), we can formulate an explicit and targeted interpolation method that takes in account all of the above.

In other words,

\begin{itemize}
\item if the wavelength $\lambda$ is much smaller than the context size $L$, we do not interpolate; 
\item if the wavelength $\lambda$ is equal to or bigger than the context size $L$, we want to only interpolate and avoid any extrapolation (unlike the previous "NTK-aware" method);
\item dimensions in-between can have a bit of both, similar to the "NTK-aware" interpolation.
\end{itemize}

As a result, it is more convenient to introduce the ratio $r = \frac{L}{\lambda}$ between the original context size $L$ and the wavelength $\lambda$. This ratio represents the number of rotations a certain RoPE dimension makes given a fixed pretrained context length $L$. In the $d$-th hidden state, the ratio $r$ depends on $d$ in the following way:

\begin{equation}
    r(d) = \dfrac{L}{\lambda_d} = \dfrac{L}{2\pi b^{\frac{2d}{|D|}}}.
\end{equation}

In order to define the boundary of the different interpolation strategies as above, we introduce two extra parameters $\alpha, \beta$. All hidden dimensions $d$ where $r(d) < \alpha$ are those where we linearly interpolate by a scale $s$ (exactly like PI, avoiding any extrapolation), and the $d$ where $r(d) > \beta$ are those where we do not interpolate at all. Define the ramp function $\gamma$ to be

\begin{align}
\gamma(r) = 
\begin{cases}
    0, &\text{if } r < \alpha\\ 1, &\text{if } r > \beta\\
    \dfrac{r - \alpha}{\beta - \alpha}, &\text{otherwise}.
\end{cases}
\end{align}

With the help of the ramp function, the "NTK-by-parts" method can be described as follows.

\begin{definition}
The "NTK-by-parts" interpolation is a modification of RoPE using Eq.~\ref{eq:modifyrope} with the following functions\footnote{The interpolation by linear ramp on $h$ may have alternatives, such as a harmonic mean over $\theta_d/s$ and $\theta_d$ converted from a linear interpolation on wavelengths. The choice of $h$ here was for the simplicity of implementation, but both would work.}.
    \begin{align}
    g(m) &= m\\
    h(\theta_d) &= 
    \mathbf{ig}(1-\gamma\big(r(d)\big)\mathbf{ig}) \frac{\theta_d}{s} + \gamma\big(r(d)\big) \theta_d.
\end{align}
\end{definition}

The values of $\alpha$ and $\beta$ should be tuned on a case-by-case basis. For example, we have found experimentally that for the Llama family of models, good values for $\alpha$ and $\beta$ are $\alpha = 1$ and $\beta = 32$.

Using the techniques described in this section, a variant of the resulting method was released under the name "NTK-by-parts" interpolation~\citep{blocntkparts}. This improved method performs better than the previous PI~\citep{chen2023extending} and "NTK-aware"~\ref{sec:ntkaware} interpolation methods, both with non-fine-tuned models and with fine-tuned models, as shown in~\citep{blocntkparts} and Section~\ref{sec:long-sequence}.

\subsection{YaRN}\label{sec:yarn}

In addition to the previous interpolation techniques, we also observe that introducing a temperature $t$ on the logits before the attention softmax has a uniform impact on perplexity regardless of the data sample and the token position over the extended context window (See Appendix~\ref{appendix:yarn}). More precisely, instead of Eq.~\ref{eq:attention}, we modify the computation of attention weights into

\begin{equation}\label{eq:attention_scaled}
\text{softmax}\left(\dfrac{{\bm{q}}_m^T{\bm{k}}_n}{t\sqrt{|D|}}\right).
\end{equation}

The reparametrization of RoPE as a set of 2D matrices has a clear benefit on the implementation of this attention scaling:  we can instead use a "length scaling" trick which scales both ${\bm{q}}_m$ and ${\bm{k}}_n$ by a constant factor $\sqrt{1/t}$ by simply scaling the complex rotary position embeddings by the same amount. With this, YaRN can effectively alter the attention mechanism without modifying its code. Furthermore, it has zero overhead during both inference and training, as rotary position embeddings are generated in advance and are reused for all forward passes. Combining it with the "NTK-by-parts" interpolation, we have the YaRN method.

\begin{definition}
    By the "YaRN method", we refer to a combination of the attention scaling in Eq.~\ref{eq:attention_scaled} and the "NTK-by-parts" interpolation introduced in Section~\ref{sec:ntkparts}.
\end{definition}

For LLaMA and Llama 2 models, we recommend the following values:

\begin{align}\label{eq:yarn_t}
    \sqrt{\frac{1}{t}} = 0.1 \ln({s}) + 1.
\end{align}

The equation above is found by fitting $\sqrt{1/t}$ at the lowest perplexity against the scale extension by various factors $s$ using the "NTK-by-parts" method (Section~\ref{sec:ntkparts}) on LLaMA 7b, 13b, 33b and 65b models without fine-tuning. We note that the same values of $t$ also apply fairly well to Llama 2 models (7b, 13b and 70b). It suggests that the property of increased entropy and the temperature constant $t$ may have certain degree of "universality" and may be generalizable across some models and training data.

The YaRN method combines all our findings and surpasses all previous methods in both fine-tuned and non-fine-tuned scenarios. Thanks to its low footprint, YaRN allows for direct compatibility with libraries that modify the attention mechanism such as Flash Attention~2~\citep{dao2023flashattention2}.

\subsection{Dynamic Scaling - "Dynamic NTK" interpolation}\label{sec:dynamic}

In a lot of use cases, multiple forward-passes are performed with varying sequence lengths from $1$ to the maximal context size. A typical example is the autoregressive generation where the sequence lengths increment by $1$ after each step. There are two ways of applying an interpolation method that uses a scale factor $s$ (including PI, "NTK-aware", "NTK-by-parts" and YaRN):

\begin{enumerate}
    \item Throughout the whole inference cycle, the embedding layer is fixed including the scale factor $s=L'/L$ where $L'$ is the fixed number of extended context size.
    \item In each forward-pass, the position embedding updates the scale factor $s=\text{max}(1, l'/L)$ where $l'$ is the sequence length of the current sequence.
\end{enumerate}

The problem of (1) is that the model may experience a performance discount at a length less than $L$ and an abrupt degradation when the sequence length is longer than $L'$. But by doing Dynamic Scaling as (2), it allows the model to gracefully degrade instead of immediately breaking when hitting the trained context limit $L'$. We call this inference-time method the Dynamic Scaling method. When it is combined with "NTK-aware" interpolation, we call it "Dynamic NTK" interpolation. It first appeared in public as a reddit post in~\citep{emozillareddit}.

One notable fact is that the "Dynamic NTK" interpolation works exceptionally well on models pretrained on $L$ without any finetuning ($L'=L$). This is supported by the experiment in Appendix~\ref{appendix:dynamic}.

Often in the repeated forward-passes, the kv-caching~\citep{kvcaching} is applied so that we can reuse the previous key-value vectors and improve the overall efficiency. We point out that in some implementations when the rotary position embeddings are cached, some care has to be taken in order to modify it for Dynamic Scaling with kv-caching. The correct implementation should cache the kv-embeddings before applying rotary position embeddings, as the RoPE of every token changes when $s$ changes.

\section{Experiments}\label{sec:experiments}

\subsection{Training}
We broadly followed the training and evaluation procedures as outlined in~\cite{chen2023extending}.

For training the 128k context window size models, we extended the Llama~2~\citep{touvron2023llama2} 7B and 13B parameter models. No changes were made to the LLaMA model architecture other than the calculation of the embedding frequencies as described in Section~\ref{sec:yarn} with $s=16$ and $s=32$.

We used a learning rate of $2 \times 10^{-5}$ with no weight decay and a linear warmup of 20 steps along with AdamW~\citep{loshchilov2018decoupled} $\beta_1=0.9$ and $\beta_2=0.95$.
For the $s=16$ model, we fine-tuned for 400 steps with global batch size $64$ using PyTorch~\citep{pytorch} Fully Sharded Data Parallelism~\citep{zhao2023pytorch} and Flash Attention~2~\citep{dao2023flashattention2} on the PG19 dataset~\citep{Rae2020Compressive} chunked into 64k segments bookended with the BOS and EOS token. For $s=32$ we followed the same procedure, but due to compute constraints, we started from the finished $s=16$ checkpoint and trained for only an additional 200 steps. Note that the $s=32$ model is also trained with 64k context data, but we show that it is able to extrapolate to a context size of 128k in Section~\ref{sec:long-sequence}.

For the ablation studies, we used the LLaMA~7B model. It has the same architecture as the newer Llama~2 models except for a shorter pretrained context window size\footnote{LLaMA models have a pretrained context size of 2k tokens, while Llama~2 models have 4k.}, which reduces compute requirements and allows for faster training and evaluations. The training procedure is similar to the 128k models, but we chunk the PG19 dataset into 32k segments instead, and train using $s=16$ for 400 steps. As shown in Figure~\ref{fig:loss32k}, YaRN converges faster compared to other interpolation techniques during training and consistently has lower loss.

\subsection{Long Sequence Language Modeling}\label{sec:long-sequence}

To evaluate the long sequence language modeling performances, we use the GovReport~\citep{huang-etal-2021-efficient} and Proof-pile~\citep{proofpile} datasets both of which contain many long sequence samples. For all evaluations, the test splits of both datasets were used exclusively. All perplexity evaluations were calculated using the sliding window method from~\citet{press2022train} with $S = 256$, which takes in account the entire documents' perplexity contribution, even if the context window of the model is shorter.

First, we select 10 random samples from Proof-pile with at least 128k tokens each and evaluate the perplexity of each of these samples when truncated at 2k steps from a sequence length of 2k tokens through 128k tokens. Table~\ref{tab:128k-comparison} shows the long sequence performance of fine-tuned Llama~2 $s=16$ and $s=32$ models. We demonstrate that YaRN is able to generalize and extrapolate to unseen context lengths and benefit from transfer learning, since the $s=32$ model was only further trained for 200 steps using the $s=16$ checkpoint with 64k data and is able to extrapolate to 128k context.

In order to further confirm the effectiveness of YaRN, we compare all four interpolation methods in Figure~\ref{fig:ppl32k} on the left and Table~\ref{tab:32k-comparison} from Appendix~\ref{appendix:ablation} as an ablation study. YaRN consistently outperforms (has lower perplexity than) other methods in both non fine-tuned and fine-tuned scenarios when using the same number of training steps. We also demonstrate that YaRN has better training efficiency compared to PI in Appendix~\ref{appendix:efficiency}. More comparisons against open models can be found in Appendix~\ref{appendix:long-sequence}.

\subsection{Passkey Retrieval}

The passkey retrieval task as defined in~\cite{mohtashami2023landmark} measures a model's ability to retrieve a simple passkey (i.e., a five-digit number) from amongst a large amount of otherwise meaningless text. For our evaluation of the fine-tuned 32k LLaMA 7B models, we performed 50 iterations of the passkey retrieval task with the passkey placed at a random location uniformly distributed across the evaluation context window on different prompt lengths ranging from 2k to 32k. YaRN achieves higher scores compared to other interpolation methods when given similar training budget, as seen in Figure~\ref{fig:ppl32k} on the right. More results and comparisons for Llama~2 models are shown in Appendix~\ref{appendix:passkey}.

\subsection{Standardized Benchmarks}

The Hugging Face Open LLM Leaderboard~\citep{openllm} compares a multitude of LLMs across a standardized set of four public benchmarks. Specifically, we use 25-shot ARC-Challenge~\citep{clark2018think}, 10-shot HellaSwag~\citep{zellers2019hellaswag}, 5-shot MMLU~\citep{hendryckstest2021}, and 0-shot TruthfulQA~\citep{lin-etal-2022-truthfulqa}.

To test the degradation of models' short context performance under context extension, we evaluated our Llama~2 and 32k LLaMA 7B models using this suite and compared it to established scores for the baselines. The results are summarized in Table~\ref{tab:open_llm} and Table~\ref{tab:open_llm_128k}. More results for Llama~2 models are shown in Appendix~\ref{appendix:standard}.

We observe that there is minimal performance degradation between the YaRN models and their respective Llama 2 baselines. Some variance is to be expected as the PG19 dataset~\citep{Rae2020Compressive} we used for fine-tuning is very different from the original pre-training datased used for LLaMA and Llama~2 models. We also observe that there was on average a 0.49\% drop in scores between the YaRN $s=16$ and $s=32$ models and can conclude that the the iterative extension from 64k to 128k results in negligible performance loss.

\subsection{Computational Efficiency}

Given that rotary position embeddings are cached during training and inference when the context window size is fixed to a preset length $L$, modifying the interpolation on rotary position embeddings incurs no additional computational or memory cost compared to previous context extension methods, which is the case for all four interpolation methods outlined in this work. YaRN converges the fastest during training compared to other methods, thus is the most computationally efficient, as shown in Table~\ref{tab:trainingtime}.

\section{Conclusion}
In conclusion, we have shown that YaRN improves upon all existing RoPE interpolation methods and can act as a drop-in replacement to PI, with no downsides and minimal implementation effort. The fine-tuned models preserve their original abilities on multiple benchmarks while being able to attend to a very large context size. Furthermore, YaRN allows efficient extrapolation with fine-tuning on shorter datasets and can take advantage of transfer learning for faster convergence, both of which are crucial under compute-constrained scenarios. Finally, we have shown the effectiveness of extrapolation with YaRN where it is able to "train short, and test long".

\section{Reproducibility}

To aid in reproducibility, we provide, as supplementary material, the entirety of of the code used to train the YaRN models in Table \ref{tab:proofpile-long-small}, as well as the evaluation code that produced Figure \ref{fig:proofpile-long-small} and Tables \ref{tab:8k-comparison}, \ref{tab:proofpile-long-small}, \ref{tab:open_llm}, \ref{tab:govreport}, and \ref{tab:passkey}. The code also contains implementations of various extension methods referenced throughout the paper. For training YaRN, we used the publicly available PG19 dataset~\citep{Rae2020Compressive} tokenized to contiguous chunks of 64k tokens.

\end{document}
